\section*{Bab \ref{ch:b1:predominant-reaction} --- Kesetimbangan Asam--Basa dan Metode Reaksi Predominan}

\begin{solution}{exo:b1:predominant-reaction:1}
Basa konjugasi: \ce{HSO4-}, \ce{CO3^{2-}}, \ce{NH3}, \ce{OH-}. Asam
konjugasi: \ce{NH4+}, \ce{H2PO4-}, \ce{H2O}, \ce{H3O+}. \omterm{def:b1:predominant-reaction:ampholyte}{Amfolit}:
\ce{HCO3-}, \ce{H2O}, \ce{HPO4^{2-}} (dan \ce{HSO4-}, basa asam sulfat
dan asam pasangan keduanya).
\end{solution}

\begin{solution}{exo:b1:predominant-reaction:2}
\ce{HF} ($K_a = \num{6.9e-4}$) $>$ \ce{CH3COOH} (\num{1.7e-5}) $>$ \ce{HClO}
(\num{2.8e-8}) $>$ \ce{NH4+} (\num{5.6e-10}).
\end{solution}

\begin{solution}{exo:b1:predominant-reaction:3}
\omterm{def:b1:predominant-reaction:strong-acid}{Asam kuat}: $\mathrm{pH} = 2.00$. \omterm{def:b1:predominant-reaction:strong-acid}{Basa kuat}: $[\ce{OH-}] = 10^{-2}$,
$\mathrm{pH} = 14.00 - 2.00 = 12.00$.
\end{solution}

\begin{solution}{exo:b1:predominant-reaction:4}
\ce{CH3COOH} di bawah pH 4.76, \ce{CH3COO-} di atasnya. Pada pH 3 bentuk
asam predominan (98\,\%); di dalam darah (pH 7.4) ion etanoat. Pada
pH 3.76 bagian asamnya $1/(1 + 10^{-1}) = 0.91$.
\end{solution}

\begin{solution}{exo:b1:predominant-reaction:5}
\omterm{def:b1:predominant-reaction:predominant-reaction}{Reaksi predominan} \ce{NH3 + H2O <=> NH4+ + OH-}, $K = 10^{-(14.00 -
9.25)} = 10^{-4.75}$. Jika hanya sedikit \ce{NH3} yang bereaksi:
$[\ce{OH-}] = \sqrt{K C} =
\qty{1.3e-3}{mol/L}$, $\mathrm{pH} = 11.13$. Pemeriksaan: $[\ce{NH4+}] = 1.3\,\%$ dari $C$, dan $\mathrm{pH} > \mathrm{p}K_a + 1$: sahih.
\end{solution}

\begin{solution}{exo:b1:predominant-reaction:6}
$\mathrm{pH} = \frac12(4.76 + 2.00) = 3.38$ (nilai eksaknya 3.39);
$\alpha = 10^{-3.38}/0.010 = 0.042$: 4\,\% terdisosiasi, di bawah
10\,\%, dan $\mathrm{pH} < \mathrm{p}K_a - 1$: sahih.
\end{solution}

\begin{solution}{exo:b1:predominant-reaction:7}
\ce{HF + OH- -> F- + H2O}, $K = 10^{14.00 - 3.16} = 10^{10.84}$: tuntas.
\omterm{def:b1:predominant-reaction:predominant-reaction}{Larutan ekuivalennya} mengandung masing-masing \qty{0.050}{mol/L} \ce{HF}
dan \ce{F-}, sebuah penyangga: $\mathrm{pH} = \mathrm{p}K_a = 3.16$.
\end{solution}

\begin{solution}{exo:b1:predominant-reaction:8}
$\mathrm{pH} = 4.76$. Asam klorida mengubah \qty{0.010}{mol} etanoat
menjadi asam: $\mathrm{pH} = 4.76 + \log(0.040/0.060) = 4.58$, perubahan
sebesar 0.18. Di dalam air murni pH-nya akan turun dari 7.00 ke 2.00.
\end{solution}

\begin{solution}{exo:b1:predominant-reaction:9}
Di dalam air kedua asam bereaksi tuntas menghasilkan \ce{H3O+}: pelarut
meratakannya. Asam etanoat adalah basa yang jauh lebih lemah daripada
air, sehingga di dalamnya kedua asam hanya bereaksi sebagian, sampai
tingkat yang berbeda, dan kekuatannya dapat dibandingkan. Asam terkuat
yang dapat ada di dalam air adalah \ce{H3O+}.
\end{solution}

\begin{solution}{exo:b1:predominant-reaction:10}
\ce{CO3^{2-} + H2O <=> HCO3- + OH-}, $K = 10^{-(14.00 - 10.33)} = 10^{-3.67}$;
$[\ce{OH-}] = \sqrt{10^{-3.67} \times 0.10} = \qty{4.6e-3}{mol/L}$
(4.6\,\% dari $C$), $\mathrm{pH} = 11.67$. Kebasaan kedua, \ce{HCO3- + H2O <=>
CO2 + H2O + OH-}, mempunyai $K = 10^{-7.63}$: kebasaan itu menghasilkan
$[\ce{CO2}] \approx
10^{-7.6}$, dapat diabaikan.
\end{solution}

\begin{solution}{exo:b1:predominant-reaction:11}
$[\mathrm{A}] = [\ce{H3O+}] = \alpha C$ dan $[\mathrm{HA}] = (1 - \alpha)C$,
sehingga $K_a = \alpha^2C/(1 - \alpha)$. Dengan $C = 10^{-3}$: $\alpha^2 +
0.0174\,\alpha - 0.0174 = 0$, $\alpha = 0.12$. Asam itu 12\,\%
terdisosiasi: hipotesis asam yang sedikit terdisosiasi ($\alpha <
10\,\%$) gagal, dan persamaan kuadratnya harus diselesaikan.
\end{solution}

\begin{solution}{exo:b1:predominant-reaction:12}
Asam klorida menetapkan $[\ce{H3O+}] = \qty{0.010}{mol/L}$: pH 2.00. Maka
\[
  [\ce{CH3COO-}] = \frac{K_a[\ce{CH3COOH}]}{[\ce{H3O+}]} = \frac{1.74 \times 10^{-5}
  \times 0.10}{0.010} = \qty{1.7e-4}{mol/L},
\]
dapat diabaikan dibandingkan dengan 0.010: ion
\ce{H3O+} dari \omterm{def:b1:predominant-reaction:strong-acid}{asam kuat} mendorong kesetimbangan \omterm{def:b1:predominant-reaction:strong-acid}{asam lemah} itu
kembali.
\end{solution}

\begin{solution}{pb:b1:predominant-reaction:1}
\textbf{1.} \ce{H3PO4}/\ce{H2PO4-}, \ce{H2PO4-}/\ce{HPO4^{2-}},
\ce{HPO4^{2-}}/\ce{PO4^{3-}}; misalnya \ce{H3PO4 <=> H2PO4- + H+}.
\textbf{2.} Batas-batasnya pada 2.15, 7.21, dan 12.34.
\textbf{3.} pH 1: \ce{H3PO4}; 5: \ce{H2PO4-}; 9: \ce{HPO4^{2-}}; 13:
\ce{PO4^{3-}}.
\textbf{4.} \ce{H2PO4-} dan \ce{HPO4^{2-}}.
\textbf{5.} Setiap proton meninggalkan ion yang muatan negatifnya lebih
besar, yang mengikat proton berikutnya lebih kuat.
\textbf{6.} \ce{H3PO4 + H2O <=> H2PO4- + H3O+}; rumus
$\frac12(\mathrm{p}K_a + \mathrm{p}C) = 1.58$ melanggar $\mathrm{pH} <
\mathrm{p}K_a - 1 = 1.15$: sekitar seperlima asamnya bereaksi.
\textbf{7.} $x^2 + K_{a1}x - K_{a1}C = 0$ dengan $K_{a1} = 10^{-2.15}$: $x =
\qty{0.0233}{mol/L}$, $\mathrm{pH} = 1.63$.
\textbf{8.} \ce{2H2PO4- <=> H3PO4 + HPO4^{2-}}, $K = 10^{2.15 - 7.21} =
10^{-5.06}$.
\textbf{9.} $\mathrm{pH} = \frac12(2.15 + 7.21) = 4.68$.
\textbf{10.} $\mathrm{pH} = \frac12(7.21 + 12.34) = 9.78$.
\textbf{11.} Tidak satu pun: 1.63, 4.68, dan 9.78; diperlukan campuran
dua di antaranya.
\textbf{12.} \ce{H3PO4 + OH- -> H2PO4- + H2O}, $K = 10^{14.00 - 2.15} =
10^{11.85}$: \omterm{def:b1:extent-q-and-k:quantitative}{kuantitatif}.
\textbf{13.} \qty{0.100}{mol/L} \ce{H2PO4-}: pH 4.68.
\textbf{14.} \qty{0.050}{mol} hidroksida berikutnya mengubah separuhnya:
\qty{0.050}{mol/L} \ce{H2PO4-} dan \ce{HPO4^{2-}}, pH 7.21.
\textbf{15.} \qty{0.100}{mol/L} \ce{HPO4^{2-}}: pH 9.78.
\textbf{16.} pH naik dari 1.63, melonjak di sekitar $n = 0.100$ (ke 4.68),
naik perlahan melewati 7.21 pada $n = 0.150$, melonjak di sekitar
$n = 0.200$ (ke 9.78), melewati 12.34 di dekat $n = 0.250$, dan mencapai
sekitar 12.6 pada $n = 0.300$ (\qty{0.100}{mol/L} \ce{PO4^3-}: $x^2/(0.100 - x) = 10^{-1.66}$, $x = 0.037$, pH 12.57).
\textbf{17.} Larutan itu mengandung asam dan basa sebuah pasangan
ber-$\mathrm{p}K_a$ 7.21 dalam jumlah yang sama
(\cref{prop:b1:predominant-reaction:buffer}).
\textbf{18.} $10^{7.20 - 7.21} = 0.977$.
\textbf{19.} $n(\ce{H2PO4-}) = 0.100 - x$, $n(\ce{HPO4^{2-}}) = x$.
\textbf{20.} $x/(0.100 - x) = 0.977$, $x = \qty{0.0494}{mol}$.
\textbf{21.} Asam itu mengubah \qty{0.0010}{mol} \ce{HPO4^{2-}}:
$\mathrm{pH} = 7.21 + \log(0.0484/0.0516) = 7.18$, perubahan sebesar
$-0.02$. Di dalam air murni: dari 7.00 ke 3.00.
\textbf{22.} Tidak: pH bergantung pada rasio jumlah-jumlahnya, yang tidak
diubah oleh pengenceran (selama \omterm{def:b1:extent-q-and-k:activity}{aktivitasnya} tetap mendekati
konsentrasinya).
\textbf{23.} $0.100 + 0.0494 = \qty{0.149}{mol}$ hidroksida:
\textbf{\qty{149}{mL}} larutan \qty{1.00}{mol/L}, yang diencerkan
sampai \qty{1.00}{L}.
\end{solution}
